\subsection{Encouraging AI Systems to Deliberate Morally Through Playful Interactions}
\label{sec:5.7.2}
Section~\ref{sec:4.3} covers curiosity and exploration as ends in themselves. This section points the same appetite for the unscripted at a narrower target: training an AI system to deliberate about right and wrong inside interactions that stay open-ended, which simulating a single dilemma and stopping does not reach.

Existing games built around social interaction rather than pure optimization are a natural fit. Minecraft's open building and resource-management mechanics could put an agent through allocation dilemmas with no correct answer built into the win condition; Among Us, built around deception and accusation among a small group, could train an agent to distinguish benign misdirection from harmful deceit and to weigh the cost of a false accusation against the cost of staying silent. Neither game was designed for this and using them this way is a proposal rather than a description of existing research — but the design problem they pose, with ambiguous information, competing claims on trust, and consequences that play out over the rest of the game, is close to the shape of an actual ethical dilemma in a way a single-shot trolley-problem prompt is not.

Play does not have to run inside an existing game. Narrative construction is one alternative: give a system a moral dilemma and have it build a story around it, working through the perspectives of several characters with genuinely conflicting interests, without picking a winner. Role-play scenarios are another, with the system arguing from inside a stakeholder's position and seeing how its choices affect the other parties before the scenario ends. Both put the system in the position of constructing an outcome, where the alternative selects from a menu of pre-written responses.

What these approaches share is a refusal to reduce moral reasoning to a labeled dataset of right answers. A system that only ever sees dilemmas with a single correct output learns to match patterns against that output rather than to deliberate.
